\documentclass[11pt,a4paper]{article}
\usepackage[T1]{fontenc}
\usepackage{lmodern,amsmath,amssymb,amsthm,mathtools,booktabs,microtype}
\usepackage[margin=27mm]{geometry}
\usepackage[colorlinks=true,linkcolor=blue,citecolor=blue,urlcolor=blue]{hyperref}
\setlength{\emergencystretch}{2em}
\newtheorem{theorem}{Theorem}[section]
\newtheorem{lemma}[theorem]{Lemma}
\newtheorem{proposition}[theorem]{Proposition}
\newtheorem{corollary}[theorem]{Corollary}
\newcommand{\Q}{\mathbb Q}\newcommand{\Z}{\mathbb Z}
\DeclareMathOperator{\Tr}{Tr}
\DeclareMathOperator{\N}{N}
\title{Admissible complement levels and prescribed native traces}
\author{}\date{16 November 2026}
\begin{document}\maketitle
\begin{abstract}
We extend the negative-odd-valuation implication for complement equations
to fixed rational levels having nonnegative even valuations outside the
source's exceptional set. For the actual conditional elliptic detector
at inputs $X=\pm1$, we construct two simultaneous rational complement
witnesses at one common slope, with the required height-dependent second
argument. The exceptional set is unchanged. Every point, difference and
trace guard, and all 128 complement guards over the two finite point
algebras, are checked exactly. Fixed positive coefficients supported on
the exceptional primes convert these witnesses into prescribed-trace
norm equations with levels $17m_i^2$. The coefficient member is $B=3$;
no sufficient defining threshold or uniform realization theorem is
asserted.
\end{abstract}

\section{Admissible fixed rational levels}

Let $S$ be a fixed set of exceptional rational primes containing $17$.
A nonzero rational level $a$ is \emph{admissible} if
\begin{equation}\label{eq:admissible}
             v_p(a)\in2\Z_{\ge0}\qquad(p\notin S).
\end{equation}
A nonzero rational class representative $d$ is allowed if
$v_p(d)\in2\Z$ for every $p\notin S$.

\begin{lemma}\label{lem:local}
If $a$ is admissible and
\[
                         q^2=d\,a(a-b)
\]
has a rational solution, then $b$ has no negative odd valuation outside
$S$.
\end{lemma}
\begin{proof}
If $v_p(b)<0$ were odd, then it would be strictly smaller than
$v_p(a)$, and
\[
 v_p\bigl(d\,a(a-b)\bigr)=v_p(d)+v_p(a)+v_p(b)
\]
would be odd. This contradicts squareness. If $a=b$, the conclusion
follows directly from~\eqref{eq:admissible}.
\end{proof}

Finitely many admissible levels and allowed classes remain fixed
coefficient data. Their finite union can be encoded by taking a product
of the separately guarded finite-algebra encodings, as
in~\cite{FiniteClasses}. This does not introduce a retained coordinate.
For arbitrary admissible levels, the norms of $a-b$ in the whole point
algebra must be established separately. No automatic nonvanishing
assertion is being made for all such levels.

\begin{lemma}\label{lem:scale}
Let $\gamma\in\Q_{>0}$ satisfy $v_p(\gamma)=0$ for every $p\notin S$.
Replacing a physical detector weight $w$ and its output $b$ by
$\gamma w,\gamma b$, while keeping the detector product unchanged,
preserves their valuations at every prime outside $S$. If the level is
also replaced by $\gamma a$, then the whole-algebra complement guard is
multiplied by the nonzero rational scalar $\gamma$.
\end{lemma}
\begin{proof}
Both assertions follow from
$v_p(\gamma)=0$ and $\gamma a-\gamma b=\gamma(a-b)$.
\end{proof}

This is a statement about fixed coefficients and local implications.
A completeness theorem for the same coefficients at all integer inputs
requires an additional argument.

\section{An identity for the actual reciprocal trace}

Put $K=\Q(\sqrt2)$, and write $W=W_0+W_1\sqrt2$ with
$W_0\N(W)\ne0$. For $w\in\Q^\times$ define
\[
 V=\Tr(W^{-1}),\qquad
 b=\frac{4w}{V^2}
   =w\left(\frac{\N(W)}{W_0}\right)^2,\qquad a=w\N(W).
\]
Then
\begin{equation}\label{eq:identity}
 2a(a-b)=
 \left(\frac{2w\N(W)W_1}{W_0}\right)^2.
\end{equation}
Indeed $W_0^2-\N(W)=2W_1^2$. The identity alone does not make $a$
admissible or fixed independently of the tuple. The next section
establishes both requirements for an explicit actual tuple.

\section{Two coupled native detector arguments}

Use the curve and points of~\cite{Direct}:
\[
 E:\ y^2=x^3-Dx,\qquad D=(3-\sqrt2)^2,\qquad
 R_0=(1+\sqrt2,2\sqrt2),\qquad P=2R_0.
\]
Thus
\[
 P=\left(\frac{57}{8}-\frac72\sqrt2,\,
          -\frac{165}{8}+\frac{417}{32}\sqrt2\right).
\]
Write $f(Q)=1/x(Q)$ and $z(Q)=-x(Q)/y(Q)$.
The source point tuple and fixed coefficient are
\begin{equation}\label{eq:tuple}
                X=1,\quad B=3,\quad A=P,\quad H=U=2P.
\end{equation}
The two comparison ratios are both one. The source scale, with exponent
four, is
\[
 a_1=(X^2+B)^4+\Tr(x(A))^4
     =\frac{10621537}{256}.
\]
Choose the common slope
\begin{equation}\label{eq:slope}
 \lambda=\frac{32003072}{10621537},\qquad
 t_0=\lambda,\quad t_1=a_1\lambda=125012.
\end{equation}
In particular, the second argument is determined by the actual root
height and the same slope.

Put
\[
 \Theta=\left\{0,-\frac12,-\frac14,-\frac25,-1,-\frac13\right\},
 \quad F(t)=\prod_{\theta\in\Theta}(t-\theta),\quad\delta=t-1,
\]
\[
 w(t)=\frac{F(t)}{4\delta^8},\qquad
 \alpha(t)=\frac{\delta^4}{F(t)}.
\]
At~\eqref{eq:tuple}, the six source tests give
\begin{equation}\label{eq:product}
 T_i=(1+\alpha(t_i)f(2P)^2)
       (1+\alpha(t_i)\sigma(f(P))^2),\qquad W_i=T_i^2,
\end{equation}
where $\sigma$ is the nonidentity automorphism of $K$. Define the two
fixed levels, physical inputs and witnesses by
\begin{equation}\label{eq:levels}
 \begin{split}
 a_i&=w(t_i)\N(W_i),\\
 b_i&=w(t_i)\left(\frac{\N(W_i)}{(W_i)_0}\right)^2,\\
 q_i&=\frac{2w(t_i)\N(W_i)(W_i)_1}{(W_i)_0}\qquad(i=0,1).
 \end{split}
\end{equation}
All constants in these formulas are specified rationally over $K$ by
\eqref{eq:tuple}--\eqref{eq:product}; they are frozen for the remainder
of the construction.

\begin{theorem}\label{thm:fixture}
For the original exceptional set $S=\{p:p\le76176\}$, both levels
$a_i$ are positive and admissible. The unchanged native inputs satisfy
\[
                 0<b_i<a_i,\qquad q_i^2=2a_i(a_i-b_i)
                 \quad(i=0,1).
\]
Every inherited point, difference and detector denominator is nonzero
at every geometric character of the point algebra, and $a_i-b_i$
is nonzero at each of its 64 characters, for both blocks.
\end{theorem}
\begin{proof}
The denominators of the two rational coefficients of $f(P)$ and $f(2P)$
have prime support $\{41,593,8273\}$. Define
\[
 \begin{split}
 P_0={}&11\cdot47\cdot53\cdot59\cdot73\cdot197\cdot227\cdot281
 \cdot653\cdot853\cdot857\cdot2203\\
 &{}\cdot8863\cdot21751\cdot31253,\\
 P_1={}&5\cdot7\cdot11\cdot13\cdot17\cdot19\cdot73\cdot137
 \cdot1697\cdot5051\cdot5483\cdot5953\cdot31253.
 \end{split}
\]
Exact substitution gives the reduced factorizations
\[
 w(t_0)=
 \frac{2^6\,3^3P_0\,10621537^2}{5^7\,7^7\,43^8\,14207^8},
 \qquad
 w(t_1)=\frac{3^3P_1}{2^2\,149^8\,839^8}.
\]
The reduced denominators of $\alpha(t_0),\alpha(t_1)$ are respectively
$2^8\,3^3P_0$ and $3^3P_1$. Thus $T_i$ is integral outside $S$.
Since $\N(W_i)=\N(T_i)^2$, every outside-$S$ valuation of $a_i$ is
nonnegative and even. No exceptional prime is added.

Both arguments exceed one, so $\alpha(t_i),w(t_i)>0$. Every factor
in~\eqref{eq:product} is positive at both real embeddings.
The exact values have $(W_i)_1\ne0$, so~\eqref{eq:identity} gives
$0<b_i<a_i$ and the displayed square equations.

For completeness of the finite check, each source point block has four
scalar roots over $K$. They independently select the two ordinate signs
at the two embeddings. The exact duplication formula
\[
       h(2Q)=\frac{x(Q)^2+D}{2y(Q)},\qquad h(2Q)^2=x(2Q)
\]
supplies the required point charts for $A,H,U$. The certificate verifies
their four scalar roots and all four sign pairs for each block.

For a character indexed by the six signs of $A,H,U$, form
$H-P$ and $U-A$ independently at the two embeddings. Their abscissas
and ordinates, and all affine difference denominators, are nonzero.
The difference points are signed copies of $P$ or $3P$.
The six source tests in the first embedding are
\[
 f(H),\ f(U),\ \sigma f(H-P),\ \sigma f(U-A),\
 \frac{z(U)}{z(H)}
 -\sigma\frac{z(U-A)}{z(H-P)},\
 \frac{z(U)}{z(H)}-1,
\]
and the second embedding uses their conjugate-slot formulas. Their
products $W_+,W_-$ have nonzero sum and product. At this character,
the physical input is
\[
                   w(t_i)\left(\frac{2W_+W_-}{W_++W_-}\right)^2.
\]
The certificate evaluates its difference from $a_i$ exactly in $K$
and verifies that it is nonzero for all $2^6$ sign choices and both
$i$. These are 128 exhaustive guard checks. The scripts described in
Section~\ref{sec:certificates} independently reconstruct these values.
\end{proof}

\section{Fixed coefficients and standard norm levels}

Let $c_i$ be the positive squarefree integer representing the rational
square class of $a_i$. The preceding factorizations give
$c_0=3\cdot5\cdot7P_0$ and $c_1=3P_1$. Set
\[
 r_i=\sqrt{a_i/c_i}\in\Q_{>0},\qquad
 k_i=\operatorname{den}(r_i),\qquad m_i=k_ir_i\in\Z_{>0},\qquad
 \gamma_i=\frac{17k_i^2}{c_i}.
\]
Admissibility shows that $c_i$ and $k_i$ are supported on $S$.
Thus $\gamma_i$ has valuation zero outside $S$. All these are fixed
constants; the exact reduced fractions are included in the certificate.

\begin{corollary}\label{cor:incidence}
Keep the actual $W_i$, reciprocal traces $V_i=\Tr(W_i^{-1})$, and
coupled arguments~\eqref{eq:slope}. Replace only the physical weights
and inputs by $w_i'=\gamma_iw(t_i)$ and $b_i'=\gamma_i b_i$.
Then the actual prescribed-trace equations
\[
 (m_iV_i)^2-2Y_i^2=\frac{4w_i'}{17},\qquad
 Y_i=2m_i[\sqrt2](W_i^{-1}),\qquad i=0,1,
\]
hold simultaneously. The levels are $17m_i^2$, and all complement
guards remain nonzero.
\end{corollary}
\begin{proof}
By construction $w_i'\N(W_i)=\gamma_i a_i=17m_i^2$.
For $W_i^{-1}=u_i+v_i\sqrt2$, the identity
$(2u_i)^2-2(2v_i)^2=4/\N(W_i)$ gives the result.
Every new complement guard is $\gamma_i$ times its previously checked
value. Lemma~\ref{lem:scale} proves the outside-$S$ valuation assertion.
\end{proof}

The same constants and witnesses work at $X=-1$, with $A=-P$,
$U=-2P$ and $H=2P$. Negating $A,U,X$ negates both ratios and both
comparison differences. Their squares, all reciprocal-abscissa tests,
and the scale are unchanged. The six-sign characters are permuted,
so the whole-algebra guard statement is unchanged as well.

The two integers $m_0,m_1$ have respectively 205 and 151 decimal digits.
Their sizes have no bearing on a quantified-variable count: they are
fixed coefficient data. This result concerns $B=3$ and inputs $\pm1$.
It does not place $B=3$ beyond the source's ineffective defining
threshold, prove the original pole-preparation conditions, or establish
completeness for the same levels and coefficients at every integer.

\section{Exact certificates}\label{sec:certificates}

The package
\[
 \texttt{computation/pole\_parameters/native\_trace\_fixture\_v47/}
\]
contains two producers, \texttt{build-fixture.py} and
\texttt{build-scaling.py}, the exact rational tables
\texttt{fixture.json} and \texttt{scaling.json}, and two independent
checkers. The first checker uses SymPy's exact $\Q(\sqrt2)$ domain;
the second uses the polynomial quotient $\Q[t]/(t^2-2)$.
Neither imports a producer or the other checker. They reconstruct the
source curve and point charts, all detector values, the factorization
certificates, both norm incidences and all 128 guards. The first also
checks the negative-input sign involution directly. The tables pin the
source file by its SHA-256 hash, and both verification receipts record
the exact table hashes.

\begin{thebibliography}{9}
\bibitem{Direct}
\emph{Thirteen Variables from Direct Elliptic Index Comparison},
project source, 10 November 2026,
\nolinkurl{papers/conditional_few_variables/direct_input13_20261110_v22.tex}.
\bibitem{FiniteClasses}
\emph{Finite square-class conditions for native detector inputs},
project source, 16 November 2026,
\nolinkurl{papers/pole_parameters/finite_complement_classes_20261116_v46.tex}.
\bibitem{UnitNorm}
\emph{A single norm determinant for the 25- and 13-variable constructions},
project source,
\nolinkurl{papers/finite_algebra/unit_norm_20261114_v28.tex}.
\end{thebibliography}
\end{document}
